\section{Objeto de análisis y método}

\subsection{Análisis inverso de Codex a partir de los archivos de sesión}

Esta entrega analiza de manera inversa los archivos de historial de sesión de Codex para revelar el diseño de su Context Engineering. Cuando el usuario envía una instrucción en Codex CLI, Codex guarda el historial completo de la conversación en un archivo JSONL local. Al analizar estos archivos, se puede reconstruir el mecanismo interno de gestión de mensajes de Codex.

\begin{importantbox}{Inyección de prompt legítima}
La «inyección de prompt» de la que se habla en esta entrega se refiere a la gestión \textbf{interna} del historial de mensajes que hacen herramientas de Agent modernas como Codex, Claude Code y OpenAI Codex: es decir, antes de que el usuario envíe la primera instrucción en la interfaz TUI, la herramienta ya insertó de antemano una gran cantidad de instrucciones del sistema. Se trata de una inyección legítima e intencional en el diseño, distinta de los ataques maliciosos de Prompt Injection.
\end{importantbox}

\subsection{Resumen de la sección}

Al analizar los archivos de sesión JSONL de Codex, se puede reconstruir por completo la forma en que se gestiona su historial de mensajes, lo que permite entender el diseño de Context Engineering de los Modern Agent.

